% Luvun otsikko. Huom, * estää luvun näkymisen sisällyluettelossa.
\chapter*{TIIVISTELMA}

Lappeenrannan–Lahden teknillinen yliopisto LUT

Oma schoolisi: LUT-kauppakorkeakoulu / LUT Energiajärjestelmät / LUT Teknis-luonnontieteellinen

Oma koulutusohjelmasi: Energiatekniikka / Kauppatieteet / Kemiantekniikka / Konetekniikka / Laskennallinen tekniikka / Sähkötekniikka / Tietotekniikka / Tuotantotalous / Ympäristötekniikka

\hspace{1em}

Tekijän nimi

\hspace{1em}

\textbf{Työn suomenkielinen nimi lihavoituna (myös mahdollinen alaotsikko)}

\hspace{1em}

Diplomityö / Kandidaatin työ

\the\year{}

% Automatically insert amount of total pages, figures, tables and appendices.
\lastpageref{LastPages} sivua,  \TotalValue{totalfigures} kuvaa, \TotalValue{totaltables} taulukkoa ja \TotalValue{appendixchapters} liitettä

\begin{flushleft}
	Examiners: Tarkastaja(t): Titteli Etunimi Sukunimi ja Titteli Etunimi Sukunimi\\
	% Adjust second examiner space.
	\hspace{19.5mm} Tarkastaja(t): Titteli Etunimi Sukunimi ja Titteli Etunimi Sukunimi
\end{flushleft}

Avainsanat: Listaa tähän avainsanat, joiden perusteella tutkielma voi haettaessa löytyä

\hspace{1em}

Tiivistelmätekstin riviväli on 1, ja se tasataan molempiin reunoihin, kuten muukin opinnäytetyön leipäteksti. Tiivistelmätekstin tunnistetietoineen on mahduttava yhdelle A4-arkille. Tiivistelmään ei voi sisällyttää salassa pidettäviä tietoja, vaan se on aina laadittava julkiseksi.

Tiivistelmä on itsenäinen esitys DI-työstä/pro gradu -tutkielmasta, ja sen tulee olla ymmärrettävissä sellaisenaan. Hyvässä tiivistelmässä virkkeet ovat täydellisiä, lyhyitä ja ytimekkäitä. Tekijän mielipiteet eivät näy, vaan hän kuvaa työtään kuin ulkopuolinen raportoija. Tiivistelmässä ei viitata yksityiskohtaisesti alkuperäistekstiin. Tutkimuksen keskeisimmät tulokset sisältyvät tiivistelmään.

Opiskelijat, joiden koulusivistyskieli ei ole suomi tai ruotsi ja jotka kirjoittavat opinnäytteensä englanniksi, kirjoittavat tiivistelmän ainoastaan englanniksi.